\section{Discussion}
\label{sec:discussion}

Read together, the two stages yield an operational recipe.
\emph{Fill with the design}: fit imputation models weighted, gate
zero-inflated components, and chain dependent targets --- the
failures of doing otherwise are invisible to headline joint-geometry
metrics and surface only in the tails and in classifier two-sample
tests, exactly where transfer-program and top-income analysis is
most exposed. \emph{Select with the targets}: treat reduction to a
simulable support as estimation, let the calibration targets drive
record retention, and refit publication weights on the selected
support rather than publishing the gated weights that selected it.
The population-view harness connects the stages: the same latent
population, viewed through surveys for evaluating fill and through
administrative aggregates for evaluating representation.

The limitations are those of each stage plus one that emerges at
the seam. Stage one's results are single-imputation; imputation
variance and multiple-imputation inference are future work, and the
tail-extension question --- how far beyond the donor support imputed
tails should extend --- is deliberately left open. Stage two's
dense-calibration baseline was still improving at its optimization
budget, so the sparse-versus-dense comparison is a
finite-optimization result on a fixed support, not a convergence
claim; concentration controls and normalized-penalty sweeps remain
to be mapped. At the seam, certification currently gates the
national and state target families while congressional-district
targets ship as diagnostics --- the assignment-only contrast in
\autoref{sec:application} is the measured cost of that gap, and
promoting district families into the certified gate is the
pipeline's next milestone. Finally, the companion deployment shows
that residual data defects propagate into named policy provisions
(a tip-income carrier undercount halves a tip-exemption estimate);
the discipline the pipeline enables is not defect-free data but
disclosed, provision-level accounting of which defects bind where.
